% Exemplo: cronograma como quadro, com células sombreadas (colortbl), que imprime bem em preto e branco
% Compilar: latexmk -pdf cronograma-sombreado.tex Testado com: abntex2, xcolor, TeX Live 2024 ou mais recente
\documentclass[12pt, a4paper, brazilian]{abntex2}
\usepackage[T1]{fontenc}
\usepackage{tabularx}
\usepackage{array}
\usepackage{multirow}
\usepackage[table]{xcolor}            % cores como black!35 e \cellcolor (carrega o colortbl)

%%BEGIN:quadropreambulo
% Float "Quadro" (a abntex2 não o define): numeração própria e lista de quadros (memoir)
\newfloat{quadro}{loq}{Quadro}
\newlistof{listofquadros}{loq}{Lista de quadros}
%%END:quadropreambulo

\begin{document}

O \autoref{qua:cronograma-cor} traz o cronograma com as células do período sombreadas.

%%BEGIN:cronogramasombreado
\newcommand{\mes}{\cellcolor{black!35}}   % célula sombreada

\begin{quadro}[htbp]
  \caption{Cronograma do projeto com células sombreadas}\label{qua:cronograma-cor}
  \small
  \renewcommand{\arraystretch}{1.5}
  \begin{tabularx}{\linewidth}{|l|*{12}{>{\centering\arraybackslash}X|}}
    \hline
    \multirow{2}{*}{\textbf{Atividade}} & \multicolumn{12}{c|}{\textbf{Mês}} \\
    \cline{2-13}
     & 1 & 2 & 3 & 4 & 5 & 6 & 7 & 8 & 9 & 10 & 11 & 12 \\
    \hline
    Revisão bibliográfica & \mes & \mes & & & & & & & & & & \\
    \hline
    Métodos & & & \mes & \mes & & & & & & & & \\
    \hline
    Testes & & & & \mes & \mes & \mes & & & & & & \\
    \hline
    Resultados & & & & & & & \mes & \mes & & & & \\
    \hline
    Conclusão & & & & & & & \mes & \mes & \mes & & & \\
    \hline
    Banca & & & & & & & & & & \mes & \mes & \mes \\
    \hline
  \end{tabularx}
  \fonte{Elaborado pelo próprio autor.}
  \nota{As células sombreadas indicam os meses de execução de cada atividade.}
\end{quadro}
%%END:cronogramasombreado

\end{document}
